\documentclass[11pt,letterpaper]{article}
\usepackage[top=1.0in,bottom=1.0in,left=1.25in,right=1.25in]{geometry}
\usepackage{microtype}
\usepackage{parskip}
\usepackage{xcolor}
\usepackage{mdframed}
\usepackage{amsmath,amssymb}
\usepackage{hyperref}
\hypersetup{colorlinks=true,urlcolor=blue,linkcolor=black,citecolor=black}
\setlength{\parindent}{0pt}
\setlength{\parskip}{7pt}

%% Style for quoted reviewer text
\mdfdefinestyle{reviewquote}{
  leftmargin=0.4cm, rightmargin=0.4cm,
  innerleftmargin=0.4cm, innerrightmargin=0.4cm,
  innertopmargin=5pt, innerbottommargin=5pt,
  backgroundcolor=gray!8,
  linecolor=gray!40, linewidth=0.5pt,
  skipabove=6pt, skipbelow=6pt
}
\newcommand{\rquote}[1]{\begin{mdframed}[style=reviewquote]\textit{#1}\end{mdframed}}

\newcommand{\sref}[1]{\textbf{[Revision:~#1]}}

\begin{document}

\textbf{Response to Reviewers}\\
\textbf{Manuscript:} PONE-D-26-05900\\
\textbf{Title:} Signal Hierarchical Petri Nets: Formal Semantics of
Hierarchical Regulatory Control of Biological Systems\\
\textbf{Date:} \today

\bigskip
\hrule
\bigskip

We thank the editor and all three reviewers for their careful and constructive
evaluation. The reviewers' comments have substantially improved the manuscript,
and we are grateful for the opportunity to revise. Below we address each point
raised, in the order: Editor, Reviewer~1 (Prof.\ Haeusler), Reviewer~2
(Dr.\ Kolč{\'a}k), Reviewer~3.

Inline revision references of the form \sref{Section~X, para.~Y} indicate
where each change can be located in the revised manuscript. A track-changes PDF
is submitted alongside this letter.

\bigskip

%% ============================================================
\section*{Editor (Dr.\ Jianhong Zhou)}
%% ============================================================

\rquote{The reviewers have raised a number of concerns that need attention,
especially the contradictory claims. In addition, please disclose if there is
any artificial intelligence (AI) tools and technologies usage following our
policy.}

\textbf{On contradictory claims.} The principal contradiction identified by
the reviewers --- that the manuscript at times claims ``first-principles''
(parameter-free) threshold computation while simultaneously acknowledging that
the parameters $\theta$ and $W_s$ are derived from experiment --- has been
corrected throughout the revised manuscript. The phrase
``first-principles prediction'' has been replaced in all ten occurrences by
``topology-driven threshold computation'' or equivalent phrasing that makes the
parameter-dependence explicit. A clarifying sentence has been added wherever
the threshold formula appears: \textit{The parameters $\theta$ and $W_s$ are
derived from experimentally measured biochemical constants; the structural claim
is that the threshold is localized to at most two quantities per signal arc,
rather than emerging from a global nonlinear fit of the whole system.}

A second contradictory claim --- that SHPN provides ``unified semantics, as
opposed to a hybrid Petri net approach'' while the formal definition is an
extension of a hybrid Petri net --- has also been corrected. The revised
Introduction and Section~4 now state that SHPN is a hybrid Petri net in
the established sense of the term (coexistent continuous and stochastic
dynamics) but extends the classical hybrid PN formulation by adding a unified
signal-type layer. \sref{Introduction, final paragraph; Sec.~4.1, Definition~2
preamble.}

\textbf{On AI tools.} The revised manuscript includes an explicit disclosure
paragraph in the Acknowledgments section, consistent with PLOS~ONE policy. The
paragraph reads: \textit{``The author used AI-assisted tools (GitHub Copilot)
as a collaborative aid for drafting, editing, and formulating arguments. All
scientific insights --- the formalism, the theorems and their proofs, the
biological interpretations, and the experimental comparisons --- originate
entirely from the author. The AI tools served to sharpen prose; they did not
generate scientific content. This manuscript was written with genuine
intellectual engagement and the author accepts full scientific responsibility
for all claims.''}

\textbf{On Journal Requirements (items 1--4).}
Style requirements are followed.
The data availability statement will be updated as requested; see our response
to Phase~4 items.
Both Table~1 and Table~2 are now referred to explicitly in the text.
\sref{Sec.~3.2, after Table~1; Sec.~4.4, after Table~2.}
Reviewer-recommended citations have been evaluated; those relevant to the
scope of the paper have been added.

\bigskip

%% ============================================================
\section*{Reviewer \#1 (Prof.\ Edward Hermann Haeusler)}
%% ============================================================

We thank Prof.\ Haeusler for his generous assessment and his expert framing of
the contribution. His comments guided the most substantive improvements in
the revision.

\subsection*{Major Concern 1 --- Scope of ``first-principles'' claim}

\rquote{While thresholds are computed structurally, the parameters $\theta(t)$
and $W_s$ are derived from experimentally measured biochemical constants. The
manuscript would benefit from clarifying that the prediction is
\emph{topology-driven} but \emph{parameter-dependent}, rather than entirely
parameter-free.}

We agree entirely, and the revision adopts the reviewer's framing wholesale.
``First-principles prediction'' has been replaced throughout by
``topology-driven threshold computation.'' The claim is now stated precisely:
the threshold is \emph{structurally localized} to the two arc-level parameters
$\theta$ and $W_s$ of the signal arc $(p_s, t) \in F_s$. Once the network
topology is fixed, the threshold $M_{\mathrm{commit}} = \theta + W_s$ is
determined by at most two measurable quantities per commitment arc, rather than
requiring a global nonlinear parameter fit of the full ODE system. The
parameters themselves are experimentally derived; the structural contribution
is the localization.
\sref{Abstract; Introduction para.~3; Corollary~1 proof; Conclusion.}

\subsection*{Major Concern 2 --- Biological generality}

\rquote{The main detailed validation focuses on a single biological system.
Although additional examples are mentioned, they are not fully developed. For a
generalist journal such as PLOS ONE, broader empirical validation or
computational benchmarking across multiple systems would strengthen the
manuscript.}

We thank the reviewer for this concern. After careful consideration, we have
chosen a different response than adding new systems in equal depth, for reasons
of scientific accuracy.

During the development of signal hierarchy theory, earlier versions of the
formalism were applied to two further biological systems:

\textbf{(i) Lambda phage lysis-lysogeny switch.} An information-theoretic
analysis of the lambda phage decision, using a predecessor formalism that
preceded the full SHPN arc-type semantics, is available as a pre-print:
Simão, E.~(2024) ``Hierarchical Preemption: A Novel Information-Theoretic
Control Mechanism in Lambda Phage Decision-Making,'' arXiv:2512.22415.

\textbf{(ii) \textit{Vibrio fischeri} quorum sensing.} An application of an
earlier 13-tuple extended Bio-PN formulation is available as a pre-print:
Simão, E.~(2025) ``Unifying Weak Independence and Signal Hierarchy Theory:\\Extended Bio-PN Formalism with Application to \textit{V.\ fischeri} Quorum
Sensing,'' arXiv:2601.00036.

Both pre-prints pre-date the signal-consumption semantics formalised in the
present paper and were deliberately not submitted for peer review until the
underlying formalism reached its current form. They are cited in the revision
(Sec.~5.3) as part of the developmental record. We do not present them as
SHPN validations because the arc-type semantics and formalism they use differ
from the definition in this paper; both are being revised under the mature
semantics before formal journal submission.

The reason for restricting quantitative validation to \textit{B.\ subtilis}
is that only this system has an independent experimental benchmark --- the
Fujita and Losick (2005) $52\%$ vs.\ $\approx 5\%$ asymmetry --- against which
the topology-driven threshold prediction can be tested without fitting. Citing
pre-print results under an earlier formalism as co-equal validation would
overstated the present contribution. We believe citing the pre-prints as
developmental context, while restricting validated claims to B.~subtilis, is
the scientifically correct position.
\sref{New Sec.~5.3 ``Developmental Scope and Pre-print Record,'' p.~40 of the revision.}

\subsection*{Major Concern 3 --- Acyclicity and feedback loops}

\rquote{The requirement that signal flow graphs be acyclic is theoretically
justified, but many biological regulatory networks exhibit feedback loops. The
manuscript addresses this through remote sensing via rate functions $\Phi$, yet
the conceptual separation between ``connected'' and ``remote'' regulation may
require further clarification.}

We thank the reviewer for raising this point; it reflects a genuine conceptual
distinction that the submitted version did not express clearly enough.

The revision clarifies this in two places. First, a dedicated paragraph in
Section~3.2 (after Table~1) explains that $\Phi$ and $G_S$ operate over
\emph{different} graph layers and answer different questions: rate functions
$\Phi$ express continuous, potentially cyclic kinetics over $G_E$; signal arcs
$F_s$ express discrete, irreversible commitment steps over $G_S$. The
acyclicity requirement applies \emph{only} to $G_S$; $G_E$ is entirely
unrestricted and may contain the feedback loops and oscillatory circuits common
in metabolic and signalling networks.

Second, the Discussion (Sec.~6) opens with a paragraph that makes this
explicit: the two feedback loops that concern the reviewer --- metabolic
recycling and regulatory oscillation --- both reside in $G_E$. What the
acyclicity of $G_S$ excludes is a closed directed path through \emph{commitment
arcs}, which would mean the controller can reset through its own commitment
channel. No such cycle has been observed in biological commitment decisions to
the authors' knowledge.

We also wish to offer a sharper biological grounding for this point. Acyclicity
of $G_S$ is not an assumption we impose --- it is a \emph{formal consequence}
of the empirical observation that committed cells do not de-differentiate
through the same biochemical channel that committed them. A directed cycle in
$G_S$ would require the existence of a firing sequence that regenerates, through
the commitment arc itself, a signal token that was consumed in the act of
deciding. This is the formal statement that a butterfly can return to its larval
state through the same metamorphic channel --- an event with no biological
counterpart. The formalism is therefore \emph{constructed to honour} that
observation, not to impose an arbitrary restriction. $G_E$ cycles remain
completely unrestricted because metabolic feedback and regulatory oscillation do
not constitute commitment reversal --- they are continuous kinetics within an
already-differentiated cell, not de-differentiation events through a committed
channel.
\sref{Sec.~3.2, after Table~1 (new paragraph); Sec.~6, opening paragraph.}

\subsection*{Major Concern 4 --- Comparison with existing PN frameworks}

\rquote{A more systematic comparison with alternative Petri net extensions
would help situate the contribution. There are approaches defined on top of
universal semantic principles, using Petri-Nets, that models the BioMolecular
realm as a Turing Machine models the symbolic computational realm, via the
Turing-Church Thesis. They do not separate signals and mass, so that a bit more
of fundamental explaining would be welcome, since they catch the Intensional
Semantics of the BioMolecular processes, without needing external parameters.}

This is a well-taken point, and we are grateful to Prof.\ Haeusler for
indicating the intensional-semantics tradition. The revision adds a dedicated
paragraph in Section~4.4 (\textit{Comparison with Classical Bio-PN}) that
situates SHPN relative to four traditions:

\begin{itemize}
\item \emph{C/E nets and occurrence nets} (Rozenberg 1998; Nielsen, Plotkin,
  Winskel 1981): these formalisms capture causal dependence and concurrency
  without introducing mass. SHPN inherits their partial-order semantics for the
  $G_S$ layer but adds the mass-transfer layer $G_E$ to handle stoichiometric
  species concentrations.
\item \emph{Process algebras} (Regev \& Shapiro 2002; Priami 2005): $\pi$-calculus
  and stochastic $\pi$-calculus model molecular interaction through channel
  communication, capturing the intensional semantics the reviewer refers to.
  These approaches are complementary: they excel at mobile processes and
  communication topology; SHPN addresses commitment thresholds as structural
  properties of a fixed network.
\item \emph{Classical hybrid PN} (Alla \& David 1998; Matsuno 2003): SHPN
  extends this tradition by adding unified signal-type semantics. The
  coexistence of continuous and stochastic dynamics is inherited; the
  $G_S$/acyclicity layer is new.
\end{itemize}

Regarding the intensional-semantics tradition specifically: we agree that
formalisms based on the Turing-Church thesis for biomolecular computation
provide a principled account of cellular processes without requiring arc-type
distinctions. Our approach is a different trade-off: SHPN sacrifices some
generality (we require explicit arc-type assignment) in exchange for a
structural theorem that makes commitment thresholds computable from the
topology. Neither approach is more fundamental; they answer different questions.
\sref{Sec.~4.4, new comparison paragraph with citations.}

\subsection*{Minor Comment --- Dual-graph diagram}

\rquote{A high-level schematic diagram early in the manuscript summarizing the
dual-graph architecture would improve accessibility.}

A TikZ figure (Fig.~1, \textit{Dual-graph architecture of SHPN}) has been
added at the close of the Introduction. It shows $G_E$ (left panel, with normal
arcs, signal places, and a feedback cycle) and $G_S$ (right panel, as a DAG
with layer depths $\lambda = 0$ and $\lambda = 1$) side by side, connected by
a bridge indicating that the signal place is shared between both graphs.
\sref{End of Introduction, new Fig.~1.}

\subsection*{Minor Comment --- Streamlining}

\rquote{Some sections could be streamlined to reduce repetition between
theoretical exposition and biological discussion.}

We have reduced repetition in Sections~3 and~5 by removing or condensing
passages that re-stated definitions already given formally. The revision is
somewhat longer than the submitted version due to the second biological example
requested by this reviewer; we have tried to offset that growth by tightening
existing prose elsewhere.

\bigskip

%% ============================================================
\section*{Reviewer \#2 (Dr.\ Juri Kolč{\'a}k)}
%% ============================================================

We thank Dr.\ Kolč{\'a}k for his detailed and precise technical critique. His
formal objections were largely correct, and addressing them has made the
manuscript considerably more rigorous. We respond to each point below.

We note that Dr.\ Kolč{\'a}k recommends rejection on grounds of formal
deficiencies. We believe the specific deficiencies he identified have been fully
repaired in the revision. We address his broader concerns --- whether $\Phi$
makes the rest of the formalism redundant, and whether SHPN introduces anything
new --- in dedicated subsections below.

\subsection*{Point 1 --- $T$ undefined in Definition~1}

\rquote{[The manuscript] immediately relies on undefined concepts ($T$ has not
been defined in Def.~1).}

Correct. The revised Definition~1 now opens with: \textit{``Let $T$ be the
set of transitions of the underlying Petri net.''} $T$ is defined before any
use. \sref{Sec.~2.1, Definition~1, first line.}

\subsection*{Point 2 --- Notation $E_s$ vs.\ $F_s$}

\rquote{[The manuscript] disrespects its own established notation ($E_s$
magically turns into $F_s$).}

The inconsistency arose from a late renaming that was not propagated everywhere.
All occurrences of $E_s$ in the submitted manuscript have been replaced with
$F_s$, including within Theorem~1 (Acyclicity Requirement) and all body text.
The notation is now consistent throughout. \sref{All occurrences; verified by
global search.}

\subsection*{Point 3 --- $W_t$ missing from the 13-tuple}

\rquote{The ``formal'' definition does not include all the components utilised
in prior text (e.g.~$W_t$).}

The test-arc weight $W_t$ has been addressed by stating the convention
explicitly within the $F_t$ bullet of the 13-tuple: \textit{``Test arc
enablement requires $M(p) \geq 1$ for all $(p,t) \in F_t$ --- the sensing
threshold is uniformly~1 and is not a free parameter of the formalism.''}
The decision not to add a new symbol $W_t$ to the tuple is deliberate: since
$W_t \equiv 1$ for all test arcs, adding it as a 14th component would introduce
a symbol that carries no modelling freedom and could mislead readers into
expecting a parameter to tune. The convention is now stated explicitly instead
of being left implicit. \sref{Sec.~4.1, 13-tuple, $F_t$ bullet.}

\subsection*{Point 4 --- Theorems 2 and 3: missing constraints}

\rquote{The ``Hierarchical Preemption'' requires extensive constraints on the
shape of the undefined set of transitions $T$, which have never been made, nor
alluded to. Essentially the same claim is made in Theorem~3, where the
``irreversibility'' part of the proof, which relies on those missing constraints,
is simply hand-waved away.}

Both theorems now carry their necessary hypotheses explicitly in their
statements (not only in the proof):

\textbf{Theorem~2} now opens: \textit{``Assuming the PreemptionCheck predicate
is enforced: transition $t_j$ is enabled only if every transition $t_i$ that
produces its signal inputs via $F_s$ paths is itself enabled.''}

\textbf{Theorem~3} now states: \textit{``Assuming no other transition
regenerates $p_s$ via $F_s$ (i.e., there is no $t' \neq t$ with
$(t', p_s) \in F_s$)''} as a hypothesis, and the proof uses ``By hypothesis''
instead of an aside.

We agree with the reviewer that unstated hypotheses are a formal deficiency.
\sref{Sec.~2.2, Theorems~2 and~3.}

\subsection*{Point 5 --- Remote sensing ($\Phi$) subsumes everything}

\rquote{In Section~3.2, ``remote sensing'' is introduced in which transitions
are assigned a function that determines whether they are enabled without any
regards to the network topology or constraints on the shape of the function
itself. This formalism is far more rich than a classical Petri net, and makes
literally all other formalisms discussed in the paper obsolete, as it subsumes
them. It is also, not novel.}

This is the most substantive technical point raised by Reviewer~2, and we
thank Dr.~Kolč{\'a}k for pressing it precisely. The submitted version did not
address it adequately. We offer a response that engages with the reviewer's own
research domain.

Dr.~Kolč{\'a}k's work on most-permissive semantics (Paulev{\'e}, Kolč{\'a}k
et al., \textit{Nature Communications} 2020; Haar \& Kolč{\'a}k, CMSB 2025)
asks: \emph{what states can a network reach under the least restrictive
interpretation?} This is the question of \emph{maximal reachability} --- the
full accessible state space of $G_E$ under the most permissive firing semantics.

$G_S$ commitment answers the \emph{dual question}: \emph{what has been
permanently ruled out?}

Each firing of a commitment transition $t$ whose input arc $(p_s, t) \in F_s$
consumes tokens from signal place $p_s$ \emph{monotonically reduces the
reachable state space of $G_E$}. The committed cell still carries the same
genome --- the uncommitted states are still topologically present in $G_E$ ---
but the signal tokens that made those states accessible have been irreversibly
consumed. Acyclicity of $G_S$ guarantees this reduction is \emph{monotone}: no
subsequent $G_S$ firing can restore the eliminated reachable states through a
committed channel. The epigenetic possibility space can only narrow, never
widen, through the commitment layer.

This gives a precise formal separation from $\Phi$: rate functions provide
\emph{kinetic expressiveness} over $G_E$, answering the most-permissive
reachability question. $G_S$ provides \emph{structural certificates of
irreversibility}, answering the dual question. A model with unconstrained $\Phi$
and no $G_S$ layer can simulate any trajectory --- but it cannot \emph{prove}
that any trajectory is permanently excluded. $G_S$ is the mechanism that makes
irreversibility a theorem rather than a simulation outcome.

Concretely: two models with identical $\Phi$ but different arc-type assignments
have different $G_S$ graphs and therefore different monotone reachability
reduction sequences. The structural theorems distinguish them; $\Phi$ alone
does not. $\Phi$ is a kinetic substrate inherited from Bio-PN; SHPN's
contribution is the $G_S$ layer and the structural theorems its topology
enables.

A second, limiting argument seals the point. Consider a Petri net in which all
regulation is encoded in $\Phi$ and \emph{no arcs exist at all}. Such a net has
no token flow: every transition fires but consumes and produces nothing, so the
marking $M$ never departs from $M_0$. $\Phi$ reads the same static values
forever. The model is a system of rate equations evaluated at fixed initial
conditions --- trivially solvable, biologically inert, structurally vacuous.
The moment a signal flow arc is added, tokens are consumed on firing, the
marking changes irreversibly, and the basin boundary $M_{\text{commit}} =
\theta + W_s$ becomes computable from the arc topology alone. $\Phi$ expresses
what a transition \emph{reads}; only $F_s$ arcs determine what it
\emph{destroys}. It is destruction --- irreversible token consumption ---
that creates the one-way barrier, and that is precisely what $\Phi$ cannot
do regardless of how rich its function class is.

We therefore regard the two frameworks as formally dual and genuinely
complementary: most-permissive semantics establishes the ceiling of what can
happen; $G_S$ commitment establishes the floor of what has been permanently
ruled out. Neither subsumes the other; they answer different questions at
different levels of the architecture.
\sref{Sec.~3.2, new paragraph after Table~1.}

\subsection*{Point 6 --- ``Unified semantics'' vs.\ ``hybrid PN'' contradiction}

\rquote{In the introduction, the author states that their model provides unified
semantics, as opposed to a hybrid Petri net approach, and then proceed to define
their model as an extension of a hybrid Petri net.}

Correct; this was a real contradiction. The revised manuscript no longer
contrasts SHPN with hybrid PN approaches. Instead, contributions point~3 in
Section~4.1 now reads: \textit{``SHPN is a hybrid Petri net in the sense of
supporting coexistent continuous and stochastic dynamics, but extends classical
hybrid PN by adding unified signal-type semantics that resolve the sub-model
separation limitation of prior approaches.''} The word ``unified'' now refers to
the arc-type semantics being unified within a single arc set, not to a contrast
with hybridness. \sref{Sec.~4.1, contributions bullet~3.}

\subsection*{Point 7 --- AI generation concern}

\rquote{The text itself is rife with poorly defined or outright meaningless
phrases with the intention of making it sound ``smart'' without actually saying
anything. A typical artifact of AI generation.}

We acknowledge the concern and have added an explicit AI disclosure to the
Acknowledgments section (also described in our Editor response above). The
scientific content of the manuscript --- the formalism, the theorems, the
biological interpretations, the experimental comparisons --- originates entirely
from the author's scientific judgement. AI tools assisted with prose editing.
We have taken the reviewer's stylistic criticism seriously and revised passages
that were identified as vague or overclaimed, in particular the ``first-
principles'' language and the remote sensing section.

\bigskip

%% ============================================================
\section*{Reviewer \#3 (anonymous)}
%% ============================================================

\rquote{The manuscript presents an ambitious and potentially valuable formalism,
Signal Hierarchical Petri Nets, intended to unify metabolic and regulatory
modelling by allowing selected biological molecules to participate
simultaneously in mass-transfer reactions and consumptive signal flow arcs.
The conceptual distinction between normal arcs, test arcs, and signal flow arcs
is interesting, and the idea of using signal-token consumption to represent
irreversible biological commitment could be a useful contribution to systems
biology and formal modelling. It is a well-written manuscript that clearly
outlines the theory and analysis, which support the conclusions drawn.}

We thank Reviewer~3 for this generous assessment and for the positive
recommendation. No specific concerns were raised.

\bigskip

%% ============================================================
\section*{Summary of all changes}
%% ============================================================

For convenience, the table below maps each reviewer point to the corresponding
change in the revised manuscript.

\begin{center}
\renewcommand{\arraystretch}{1.3}
\begin{tabular}{p{4cm}p{4.5cm}p{5cm}}
\hline
\textbf{Reviewer point} & \textbf{Change} & \textbf{Location} \\
\hline
Editor: contradictory claims & Replaced all ``first-principles'' language; corrected hybrid-PN framing & Throughout; Intro; Sec.~4.1 \\
Editor: AI disclosure & Acknowledgments paragraph added & Sec.~Acknowledgments \\
Editor: Table refs & Added in-text \texttt{\textbackslash ref} calls & Sec.~3.2; Sec.~4.4 \\
R1-C1: ``first-principles'' & Replaced with ``topology-driven'' & All 10 occurrences \\
R1-C2: second biological system & Pre-print developmental record cited; B.~subtilis remains sole validated system & Sec.~5.3 (new) \\
R1-C3: acyclicity / feedback & New paragraph + Discussion opening & Sec.~3.2; Sec.~6 \\
R1-C4: PN comparison & New comparison paragraph & Sec.~4.4 \\
R1-minor: dual-graph diagram & Fig.~1 added & End of Introduction \\
R2-P1: $T$ undefined & ``Let $T$ be...'' added & Def.~1 \\
R2-P2: $E_s \to F_s$ & Global replacement & All occurrences \\
R2-P3: $W_t$ missing & Convention stated in $F_t$ bullet & Sec.~4.1 \\
R2-P4: Thm.~2/3 missing hypotheses & Hypotheses moved into statement & Theorems~2 and~3 \\
R2-P5: $\Phi$ subsumes $G_S$? & Clarifying paragraph & Sec.~3.2 \\
R2-P6: hybrid-PN contradiction & Contributions bullet~3 rewritten & Sec.~4.1 \\
\hline
\end{tabular}
\end{center}

%% ============================================================
\section*{Working note: $\Gamma = (K, n, \varepsilon)$ and $\theta_{\mathrm{eff}}$ --- epistemic status (draft, not submitted)}
%% ============================================================

This note records the two-level epistemic structure of the ``not fitted'' claim,
for integration into a future revision round if a reviewer raises the question.

\textbf{What $\Gamma = (K, n, \varepsilon)$ is.}
Each signal flow arc carries a kinetic constraint tuple with three components:
$K$ (Michaelis constant, same units as marking), $n$ (Hill coefficient), and
$\varepsilon \in [0,1)$ (suppression floor fraction).
The effective basin floor is
$\theta_{\mathrm{eff}} = K \cdot (\varepsilon / (1 - \varepsilon))^{1/n}$.
When $\varepsilon = 0$ (default), $\theta_{\mathrm{eff}} = 0$.

\textbf{Two-level epistemic structure.}
\begin{enumerate}
\item \emph{Cross-domain grounding:} $\Gamma$ IS fitted --- to molecular kinetics
  data from BRENDA and published biochemical measurements. These parameters
  describe enzyme--substrate binding at the molecular level, a different
  experimental domain from the sporulation fraction being predicted.
\item \emph{Analytical derivation:} $\theta_{\mathrm{eff}}$ and $M_{\mathrm{commit}}$
  are NOT fitted to any cellular-level observation. The formula is applied once,
  with no adjustment targeting the Fujita fractions.
\end{enumerate}

\textbf{Precise claim:} Molecular-scale kinetic constants, sourced independently from BRENDA,
predict a cell-population-scale commitment threshold without any intermediate fitting step
targeting the predicted outcome. This is \emph{cross-domain grounding}, which is stronger
than ``fitted'' (no free parameter targeting the outcome) but not ``parameter-free''
(enzyme kinetics have been measured by others).

\textbf{Current manuscript state (2026-06-23):}
\begin{itemize}
\item $\Gamma = (K, n, \varepsilon)$ appears 5 times but was never formally defined.
\item $\theta_{\mathrm{eff}}$ had 0 occurrences.
\item The formula $\theta_{\mathrm{eff}} = K \cdot (\varepsilon/(1-\varepsilon))^{1/n}$ was absent.
\item Remark added to \S6.1.3 (Corollary: Threshold Predictability) in this revision session.
\end{itemize}

\end{document}
